\subsection{Scenario I: Stationary Identity}
\label{sec:eval:s1}
Scenario~I freezes every node
($v_{\mathrm{cmd}}=v_{\mathrm{drift}}=0$)
for $T_{\mathrm{sim}}=1600\,\mathrm{s}$ and $N_{\mathrm{gen}}=2160$
at $N=40$.
E-LQTP is run in the identity configuration (trigger off,
$\varepsilon_0$ huge, EKF not ticked) so that any departure from
LQARP would be a code defect rather than a mobility effect
\cite{liu2026lqarp}.
Table~\ref{tab:s1_identity} reports the five-seed means.
LQARP and identity E-LQTP are bit-identical on every seed: PDR
$1.000$, hops $1.530\pm 0.068$, delay $3.078\pm 0.071\,\mathrm{s}$,
$E_{\mathrm{bit}}=8.20\pm 0.37\times 10^{-3}\,\mathrm{J/bit}$,
$O_{\mathrm{ctrl}}=0$, VOID $=0$.
Fixed-Periodic records $442.03\,\mathrm{byte/s}$ on this frozen
field.
A nominal budget of one $400$-bit HELLO from each of $41$ nodes
every $5\,\mathrm{s}$ is $410\,\mathrm{byte/s}$.
The recorded stationary rate exceeds that budget.
The mobile comparison in Section~\ref{sec:eval:metrics} uses the
exact grid $413.42\,\mathrm{byte/s}$ ($41$ transmitters $\times$
$121$ instants on $[0,600]\,\mathrm{s}$).
A timer still spends that control traffic on a field that cannot
change, the same class of waste a timer-rebuilt STG or a periodic
QTAR HELLO incurs on a frozen mesh
\cite{xu2025stgr,nandyala2023qtar}.
RLOR-inspired always takes the one-hop sink path
($\mathrm{hops}=1.000$) because it drops $N_{\mathrm{depth}}$ and
forwards $\arg\max P_{\mathrm{succ}}$; that is the opportunistic
construction surveyed in \cite{jiang2023orsurvey}, not a routing
win claimed against LQARP.
PDR saturates for all four protocols, which is the intended ceiling
of a stationary, range-rich column, not a primary comparison.

% Scenario I stationary identity. Single-column: four short rows.

\begin{table}[!t]
\caption{Scenario~I identity ($N=40$, $T_{\mathrm{sim}}=1600\,\mathrm{s}$,
five seeds). E-LQTP trigger off; PDR $=1.000$, VOID $=0$.}
\label{tab:s1_identity}
\centering
\scriptsize
\resizebox{\columnwidth}{!}{%
\begin{tabular}{lcccc}
\hline
Protocol & Hops & $T_{\mathrm{e2e}}$ (s) & $E_{\mathrm{bit}}$ & $O_{\mathrm{ctrl}}$ \\
\hline
LQARP & $1.530\pm 0.068$ & $3.078\pm 0.071$ & $8.20\pm 0.37\times 10^{-3}$ & $0$ \\
E-LQTP & $1.530\pm 0.068$ & $3.078\pm 0.071$ & $8.20\pm 0.37\times 10^{-3}$ & $0$ \\
RLOR-insp. & $1.000\pm 0$ & $1.887\pm 0.108$ & $5.36\times 10^{-3}$ & $0$ \\
Periodic & $1.546\pm 0.039$ & $3.119\pm 0.188$ & $8.29\pm 0.21\times 10^{-3}$ & $442.03$ \\
\hline
\end{tabular}}
\end{table}

